\documentclass[../../main.tex]{subfiles}
\begin{document}

{\bf [解1]}

\paragraph{(1)}
数学的帰納法で示す．$n=1$のときは$p_1(x)=x$，$q_1(x)=1$とすればよい．

$n=k$での成立を仮定する（$k\in\mathbb{N}$）．加法定理より，
\begin{align}
  \sin(k+1)\theta
  &=\sin k\theta\cos\theta+\cos k\theta\sin\theta\\
  &=p_k(\tan\theta)\cos^{k+1}\theta
  +q_k(\tan\theta)\cos^k\theta\sin\theta\\
  &=\bigl\{p_k(\tan\theta)+\tan\theta\,q_k(\tan\theta)\bigr\}
  \cos^{k+1}\theta,\\
  \cos(k+1)\theta
  &=\cos k\theta\cos\theta-\sin k\theta\sin\theta\\
  &=\bigl\{q_k(\tan\theta)-\tan\theta\,p_k(\tan\theta)\bigr\}
  \cos^{k+1}\theta.
\end{align}
したがって，
\begin{align}
  p_{k+1}(x)&=p_k(x)+xq_k(x),\nonumber\\
  q_{k+1}(x)&=q_k(x)-xp_k(x)\label{eq:ut1991-4-recurrence}
\end{align}
とすれば$n=k+1$でも成立する．よって帰納法により題意は示された．

\paragraph{(2)}
$k=1$として，$p_2(x)=2x$，$q_2(x)=1-x^2$だから，
\begin{align}
  p_2'(x)=2=2q_1(x),\qquad q_2'(x)=-2x=-2p_1(x)
\end{align}
となり，$n=2$では題意が成立する．$n=k$での成立を仮定する（$k\geqq2$）．\eqref{eq:ut1991-4-recurrence}より，
\begin{align}
  p_{k+1}'(x)
  &=p_k'(x)+q_k(x)+xq_k'(x)\\
  &=k\bigl(q_{k-1}(x)-xp_{k-1}(x)\bigr)+q_k(x)\\
  &=(k+1)q_k(x),\\
  q_{k+1}'(x)
  &=q_k'(x)-p_k(x)-xp_k'(x)\\
  &=-k\bigl(p_{k-1}(x)+xq_{k-1}(x)\bigr)-p_k(x)\\
  &=-(k+1)p_k(x).
\end{align}
したがって$n=k+1$でも成立する．よって帰納法により題意は示された．

\bigskip\noindent
{\bf [解2]}

\paragraph{(2)}
$h_n(x)=q_n(x)+ip_n(x)$とおく．$h_1(x)=1+ix$であり，\eqref{eq:ut1991-4-recurrence}から
\begin{align}
  h_{n+1}(x)
  &=\bigl(q_n(x)-xp_n(x)\bigr)
  +i\bigl(p_n(x)+xq_n(x)\bigr)\\
  &=(1+ix)h_n(x).
\end{align}
くり返し用いて，$h_n(x)=(1+ix)^n$を得る．したがって，
\begin{align}
  p_n(x)&=\operatorname{Im}h_n(x)
  =\dfrac{(1+ix)^n-(1-ix)^n}{2i},\\
  q_n(x)&=\operatorname{Re}h_n(x)
  =\dfrac{(1+ix)^n+(1-ix)^n}{2}.
\end{align}
両辺を微分して，
\begin{align}
  p_n'(x)
  &=\dfrac{n}{2}\bigl\{(1+ix)^{n-1}+(1-ix)^{n-1}\bigr\}
  =nq_{n-1}(x),\\
  q_n'(x)
  &=\dfrac{ni}{2}\bigl\{(1+ix)^{n-1}-(1-ix)^{n-1}\bigr\}
  =-np_{n-1}(x).
\end{align}
を得る．

\end{document}
